% Clase 19 --- El campo eléctrico inducido de una región con B(t) (§3.2).
% "Círculos centrados acá... el campo eléctrico lo voy a dibujar para allá, pero
%  podría ser para el otro lado, en función de que B esté creciendo o
%  decreciendo."
% Ojo con el sentido: la figura NO puede comprometerse, porque depende del signo
% de dB/dt. Se dibuja un sentido y se aclara explícitamente que el otro caso lo
% invierte --- eso es más honesto que elegir uno y callarlo.
\begin{center}
\begin{tikzpicture}[scale=1]

  % =============== (a) las líneas de E ===============
  \begin{scope}
    \fill[figblue!8] (0,0) circle (2.1);
    \draw[figgray, line width=0.8pt] (0,0) circle (2.1);
    \foreach \ang in {0,60,120,180,240,300} {
      \node[figblue, scale=0.7] at (\ang:1.55) {$\otimes$};
      \node[figblue, scale=0.7] at ({\ang+30}:0.8) {$\otimes$};
    }
    \node[figblue, scale=0.7] at (0,0) {$\otimes$};
    \node[figlblaux, anchor=south west] at (1.5,2.15) {$\vec B(t)$ uniforme};

    % líneas de campo eléctrico
    \foreach \rr in {0.65,1.25,1.85} {
      \draw[figred, line width=0.9pt] (0,0) circle (\rr);
      \draw[figvecres, line width=0.9pt] ({95}:\rr) arc (95:65:\rr);
    }
    \node[figlbl, figred, anchor=west] at (0:1.9) {$\vec E$};

    % la curva amperiana
    \draw[figamber, dashed, line width=1.2pt] (0,0) circle (1.55);
    \draw[figvecaux, figamber] (0,0) -- (215:1.55);
    \node[figlbl, figamber, anchor=north] at (215:0.85) {$r$};

    \node[figlblaux, anchor=north, align=center] at (0,-2.35)
      {el sentido depende del signo de $dB/dt$};
    \node[figlbl, anchor=north, align=center] at (0,-2.95)
      {(a) $E\,(2\pi r) = -\dfrac{d\Phi_B}{dt}$, \; $\Phi_B = B\pi r^2$};
  \end{scope}

  % =============== (b) el resultado ===============
  \begin{scope}[xshift=6.4cm, yshift=-0.7cm]
    \begin{axis}[
      fisfig, width=5.8cm, height=4.4cm,
      xlabel={$r$}, ylabel={$|E|$},
      xmin=0, xmax=2.6, ymin=0, ymax=1.35,
      xtick={2}, xticklabels={$R$}, ytick=\empty,
    ]
      \addplot[figred, samples=2, domain=0:2] {0.55*x};
      \draw[figaux] (axis cs:2,0) -- (axis cs:2,1.1);
      \node[figlblaux, figred, anchor=north west] at (axis cs:0.1,1.28)
        {$|E| = \dfrac{r}{2}\left|\dfrac{dB}{dt}\right|$};
    \end{axis}
    \node[figlbl, anchor=north, align=center] at (2.9,-1.4)
      {(b) crece linealmente con $r$};
  \end{scope}

\end{tikzpicture}\\[0.5em]
{\small\itshape El ejemplo es el análogo eléctrico del hilo infinito, y se
resuelve con el mismo argumento de simetría, sólo que ahora la ley que se usa
es la de Faraday. Como no hay cargas y la situación es invariante ante
rotaciones alrededor del centro, las líneas de $\vec E$ tienen que ser
circunferencias concéntricas con $|\vec E|$ constante sobre cada una ---y
cerradas, que es la marca de un campo que no deriva de un potencial---. Entonces
$\oint \vec E\cdot d\vec l = E\,(2\pi r)$, el flujo encerrado es $B\pi r^2$ y
queda $|E| = \tfrac{r}{2}|dB/dt|$: nulo en el centro y creciendo linealmente
hacia afuera. El sentido de recorrido es lo único que la figura no puede
decidir, porque depende de si $B$ está creciendo o decreciendo; dibujarlo sin
aclararlo sería tomar partido en algo que el enunciado no fija. Conviene retener
la lectura conceptual: un campo magnético variable \emph{crea} campo eléctrico
en el vacío, sin cargas de por medio.}
\end{center}
